\documentclass[11pt]{article}
\usepackage[margin=1in]{geometry}
\usepackage[T1]{fontenc}
\usepackage{lmodern,amsmath,amssymb,amsthm,microtype}
\usepackage[colorlinks=true,linkcolor=blue,citecolor=blue,urlcolor=blue]{hyperref}
\hypersetup{pdftitle={A three-point counterexample to a random lattice discrepancy comparison},
pdfauthor={Alec Kriebel},pdfsubject={A negative answer to Novak's expected discrepancy moment comparison},
pdfkeywords={Lp discrepancy, random lattice, random rotation, fourth moment, counterexample}}
\newtheorem{theorem}{Theorem}
\newtheorem{proposition}[theorem]{Proposition}
\newcommand{\E}{\mathbb E}
\newcommand{\Prob}{\mathbb P}
\newcommand{\disc}{\operatorname{disc}}
\newcommand{\iid}{\mathrm{iid}}
\title{A three-point counterexample to a random lattice discrepancy comparison}
\author{Alec Kriebel\\
\small\href{https://orcid.org/0009-0001-9320-500X}{ORCID: 0009-0001-9320-500X}}
\date{17 October 2026 \quad Version 1.0.0}
\begin{document}
\maketitle
\begin{abstract}
For three points $\{\Delta\},\{Z+\Delta\},\{2Z+\Delta\}$ with independent
continuous uniform generator $Z$ and shift $\Delta$, the expected normalized
anchored fourth discrepancy moment is $19/1620$. The corresponding
independent-point moment is $4/405=16/1620$. This gives a negative answer
to the universal expected-moment comparison posed by Novak, reporting
joint work with Hinrichs, in 2004. We give an elementary inclusion-area
proof and a separate empty-tail argument giving reversal for every
point count $n\ge4$ in dimension one. Three is the smallest possible
counterexample point count. Ferudun's preprint of 14 October 2026 also
gives the same exact counterexample and broader results; this note
records the proof in the source repository PR dated 10 October.
First priority is not claimed.
\end{abstract}

\section{The comparison and its recorded scope}
For indexed points $t_0,\ldots,t_{n-1}\in[0,1)^d$, define
\[
 \disc_p(t)^p
 =\int_{[0,1]^d}\left|
       \frac1n\sum_{j=0}^{n-1}\mathbf1_{[0,x)}(t_j)
                    -\prod_{r=1}^d x_r\right|^p\,dx .
\]
Novak, reporting joint work with Hinrichs, asks whether, for even $p$,
the expectation of this quantity for
\[
 M_n^{Z,\Delta}=(\{jZ+\Delta\})_{j=0}^{n-1},
 \qquad Z,\Delta\ \hbox{independent uniform on }[0,1)^d,
\]
is at most the expectation for $n$ independent uniform points
\cite[p.~698]{Novak}. Fractional parts are coordinatewise. This is the
full continuous generator-and-shift average, with no selection of good
generators. A single reversed triple $(n,d,p)$ disproves the universal
comparison.

\begin{theorem}\label{thm:main}
In dimension one,
\[
 \E\disc_4(M_3^{Z,\Delta})^4
 =\frac{19}{1620}
 >\frac{4}{405}
 =\E\disc_4(M_3^\iid)^4 .
\]
The excess is $1/540$ and the ratio is $19/16$.
For $n\le2$ the complete lattice and independent-point ensemble laws agree
in every dimension. Thus three is the smallest possible point count
for a counterexample.
\end{theorem}

\paragraph{Related result and chronology.}
Ferudun's unrefereed preprint \cite[Theorems~1.2--1.3]{Ferudun},
released on 14 October 2026, contains the same three-point values,
the same empty-progression mechanism, and broader results in dimension
and exponent. It is an explicit overlapping answer and is credited here.
The source proof of the present note is in repository
PR~1135, whose GitHub creation record is dated 10 October 2026
\cite{Repository}; the present note is a new release on 17 October.
These records do not certify historical public visibility of the repository,
first discovery, or first publication. No inference about either author's
awareness or dependence is made. The bounded literature audit is supplied
with the verification package; it identified no earlier full answer but
does not certify exhaustive priority. The elementary proof below is
self-contained.

\section{The three-point calculation}
Set $U=t_0=\Delta$ and $V=t_1=\{Z+\Delta\}$. These are independent uniform
variables, since translation preserves uniform circle measure, and
$t_2=\{2V-U\}$. More generally, every two distinct indexed points are
independent uniforms: conditional on $Z$, $t_j$ is uniform, so $t_j$
is independent of $Z$; multiplication by the nonzero integer $k-j$
preserves circle measure, and
$t_k=\{t_j+(k-j)Z\}$. This argument also applies coordinatewise.

For $0\le x\le1/2$, put
\[
 q(x)=\Prob(t_0<x,\ t_1<x,\ t_2<x).
\]
On $U,V\in[0,x)$, one has $-x<2V-U<2x$. A negative value wraps to
at least $1-x\ge x$; a nonnegative value does not wrap through one,
outside null boundaries. Therefore the third inclusion is equivalent to
$0\le2V-U<x$. For each $U\in[0,x]$, the allowed interval
\[
 U/2\le V<(U+x)/2
\]
lies in $[0,x]$ and has length $x/2$. Its area is
\[
 q(x)=x^2/2,\qquad 0\le x\le1/2.
\]

Let $N_x$ count the three indexed points in $[0,x)$, and let
$P_k(x)=\Prob(N_x=k)$. Pairwise independence gives
$\E N_x=3x$ and $\E\binom{N_x}{2}=3x^2$. Together with $P_3=q$ and
total probability one, this determines
\[
 (P_0,P_1,P_2,P_3)
 =(1-3x+3x^2-q,\ 3x-6x^2+3q,\ 3x^2-3q,\ q).
\]
Consequently, before substituting $q$,
\begin{equation}\label{eq:general}
 F(x):=\E(N_x/3-x)^4
 =x^4-\frac23x^3+\frac2{27}x^2+\frac{x}{27}
                         +\frac49(1-2x)q(x).
\end{equation}
On the lower half interval this becomes
\[
 F(x)=x^4-\frac{10}{9}x^3+\frac8{27}x^2+\frac{x}{27}.
\]
The parameter transformation $(Z,\Delta)\mapsto(-Z,-\Delta)$ modulo one
preserves their product uniform law and reflects every point. Outside
null endpoint events, it exchanges $N_x$ with $3-N_{1-x}$.
Thus $F(1-x)=F(x)$. Bounded measurability permits Tonelli's theorem, and
\[
 \E\disc_4(M_3^{Z,\Delta})^4
 =2\int_0^{1/2}F(x)\,dx
 =2\left(\frac1{160}-\frac5{288}+\frac1{81}+\frac1{216}\right)
 =\frac{19}{1620}.
\]

For an independent triple, the same identity \eqref{eq:general} holds
with $q_\iid(x)=x^3$. On $[0,1/2]$ its difference from the lattice
integrand is therefore
\[
 F(x)-F_\iid(x)
 =\frac49(1-2x)(x^2/2-x^3)
 =\frac29x^2(1-2x)^2 .
\]
Reflection applies to both ensembles, so the strict integrated excess is
\[
 2\int_0^{1/2}\frac29x^2(1-2x)^2\,dx=\frac1{540}.
\]
This also verifies the sign without subtracting approximate values.

For completeness, the iid fourth moment at every point count has a short
formula. For independent Bernoulli$(x)$ variables $B_i$, set $Y_i=B_i-x$
and $a=x(1-x)$. Then $\E Y_i^2=a$ and $\E Y_i^4=a-3a^2$.
Only the single-index and paired-index terms survive in the fourth
power of their centered sum. Hence
\[
 \E\left(\frac1n\sum_iY_i\right)^4
 =\frac{n(a-3a^2)+3n(n-1)a^2}{n^4}.
\]
Integrating $a$ and $a^2$ gives $1/6$ and $1/30$, respectively, yielding
\begin{equation}\label{eq:iid}
 \E\disc_4(M_n^\iid)^4=\frac{3n-1}{30n^3}.
\end{equation}
At $n=3$ this is $4/405$. The whole-law claim for $n\le2$ follows
from the uniform marginal and the independent pair already proved,
coordinate by coordinate. This completes Theorem~\ref{thm:main}.

\section{A second counterexample from empty tails}
\begin{proposition}\label{prop:tails}
For $n\ge3$ in dimension one,
\[
 \E\disc_4(M_n^{Z,\Delta})^4\ge\frac2{105(n-1)}.
\]
For every integer $n\ge4$ this lower bound exceeds the iid fourth moment.
\end{proposition}
\begin{proof}
On the open triangle
\[
 0<Z<1/(n-1),\qquad 0<\Delta<1-(n-1)Z,
\]
all points lie without wraparound between $\Delta$ and
$\Delta+(n-1)Z$. Below the first point the empirical count is zero;
above the last it is $n$. The two outer intervals alone contribute
\[
 \frac{\Delta^5+(1-\Delta-(n-1)Z)^5}{5}
\]
to the fourth discrepancy moment. Integrating this contribution against
the original product density one, first in $\Delta$ and then in $Z$, gives
\[
 \int_0^{1/(n-1)}\frac{(1-(n-1)Z)^6}{15}\,dZ
 =\frac1{105(n-1)} .
\]
The transformation $(Z,\Delta)\mapsto(1-Z,1-\Delta)$ gives a descending
no-wrap triangle with the same contribution. Their interiors are disjoint
for $n\ge3$, since their generator intervals are disjoint; boundaries
have area zero. Everything omitted is nonnegative, giving the bound.

At $n=4$ it gives $2/315>11/1920$ with margin $5/8064$.
In general, by \eqref{eq:iid}, the comparison reduces to
$4n^3-21n^2+28n-7>0$. Writing $n=4+t$ makes its left side
$25+52t+27t^2+4t^3>0$ for $t\ge0$.
\end{proof}

The lower bound alone is insufficient at $n=3$, where
$1/105<4/405$. The exact three-point calculation is separate.
The rare short progressions explain how a fourth-moment reversal can
coexist with pairwise independence.

\section{Consistency and verification}
For any $n,d$, an anchored box of volume $v=\prod_r x_r$ has pairwise
independent indicator variables in the lattice ensemble. Its normalized
count variance is $v(1-v)/n$, as for iid points. Thus
\[
 \E\disc_2(M_n^{Z,\Delta})^2
 =\frac{2^{-d}-3^{-d}}{n}
 =\E\disc_2(M_n^\iid)^2.
\]
Indexing and the weight $1/n$ are retained even at collisions. Parameters
causing collisions or tested endpoint equalities are null sets; they
are not conditioned away. Closed parameter cubes and half-open cubes
give the same expectations.

The result concerns $\E\disc_p^p$. It does not order $\E\disc_p$,
settle the separate existence question for good lattice generators,
or refute discrepancy estimates restricted to selected generators.

The source archive includes the unchanged source-PR documents, three new
independent mathematical audits, current priority and chronology reviews,
and exact rational standard-library verification programs. Besides the
inclusion calculation, a separate complete affine-cell integration gives
$19/1620$ for three points and the supplementary value $139/17280$
for four points. The latter value is not needed by either analytic
counterexample. Ordinary and optimized Python runs agree. The archive
documents source versions and inaccessible literature; successful checks
do not establish exhaustive priority or formal proof certification.

\paragraph{AI assistance and review status.}
AI tools were used extensively in proof development, exact computation,
adversarial verification, literature research and manuscript preparation.
This note and its supporting materials are unrefereed. Independent AI
adversarial reviews are not external human peer review, journal acceptance
or proof-assistant certification.

\begin{thebibliography}{9}
\bibitem{Novak}
E.~Novak (joint work with A.~Hinrichs),
\emph{New bounds for the star discrepancy},
in Discrepancy Theory and Its Applications, Oberwolfach Report~13/2004,
pp.~696--699.
\href{https://doi.org/10.4171/OWR/2004/13}{doi:10.4171/OWR/2004/13}.

\bibitem{Ferudun}
A.~Ferudun,
\emph{A Negative Answer to a Question of Novak on Randomly Shifted Lattices
and the Expected p-th Power of the $L_p$ Discrepancy},
EulerSolve Research Papers, OWR-723-005, version~1.0,
14~October 2026, unrefereed preprint.
\href{https://doi.org/10.5281/zenodo.23372668}{doi:10.5281/zenodo.23372668}.

\bibitem{Repository}
A.~Kriebel, source proof and audit for problem~30000052,
\href{https://github.com/AlecKriebel/Math/pull/1135}{Math repository PR~1135},
GitHub creation record 10~October 2026, reviewed head
\texttt{e809a9c40f4ccdaced6f45f10b1dc8b7725a2119}.
Exact source identities and dated API receipts accompany this note.
\end{thebibliography}
\end{document}
